\documentclass[10pt,a4paper]{amsart}
\usepackage[margin=19mm]{geometry}
\usepackage{amsmath,amssymb,mathtools}
\usepackage[scr=rsfs,cal=euler]{mathalfa}
\usepackage{xfrac}
\usepackage[unicode,hidelinks,bookmarks=false]{hyperref}
\newcommand{\ie}{i.e.\ }
\newcommand{\cf}{cf.\ }
\newcommand{\resp}{respectively\ }
\newcommand{\Subsection}{\subsection}
\providecommand{\nobreakdash}{\nobreak\hbox{-}}
\setlength{\emergencystretch}{2em}
\multlinegap0pt
\usepackage{mathtools}
\usepackage{soul}
\usepackage{tikz}
\usepackage{pgfplots}
\usepackage{pgfplotstable}
\usepackage{mathrsfs}
\usepackage{array}
\usepackage{pgffor}
\usepackage{float}
\renewcommand{\baselinestretch}{1.4}
\usepackage{multirow}
\newtheorem{theorem}{Theorem}[section]
\newtheorem{algorithm}{Algorithm}[section]
\newtheorem{lemma}{Lemma}[section]
\newtheorem{proposition}{Proposition}[section]
\newtheorem{corollary}{Corollary}[section]
\newtheorem{remark}{Remark}[section]
\newtheorem{definition}{Definition}[section]
\newtheorem{example}{Example}[section]
\newtheorem{property}{Property}[section]
\numberwithin{equation}{section}
\providecommand{\Real}{\mathop{\rm Re}\nolimits}
\providecommand{\Imag}{\mathop{\rm Im}\nolimits}
\providecommand{\Res}{\mathop{\rm Res}}
\DeclareMathOperator*{\esssup}{ess\,sup}
\definecolor{pomcol}{rgb}{1,0,0}
\definecolor{dbcol}{rgb}{0,0,0.8}
\newcommand{\pom}[1]{{\color{pomcol}{#1}}}
\newcommand{\ppmnote}[1]{\footnote{{POM: #1}}}
\newcommand{\pomdel}[1]{{\st{#1}}}
\newcommand{\db}[1]{{\color{dbcol}{#1}}}
\newcommand{\dbnote}[1]{\footnote{\db{DB: #1}}}
\newcommand{\dbdel}[1]{\db{\st{#1}}}
\providecommand{\Real}{\mathop{\rm Re}\nolimits}%
\providecommand{\Imag}{\mathop{\rm Im}\nolimits}%
\providecommand{\Res}{\mathop{\rm Res}}
\begin{document}
\title[Fractal-Fractional HIV Dynamics with Mittag-Leffler Kernel: ]{Fractal-Fractional HIV Dynamics with Mittag-Leffler Kernel: Analysis, Stability, and Numerical Simulations}
\author{Niaz Ali Shah, Samad Noeiaghdam}
\date{}
\maketitle
\noindent\emph{Source and licence.} Fractal-Fractional HIV Dynamics with Mittag-Leffler Kernel: Analysis, Stability, and Numerical Simulations. Version arXiv:2607.00061v1 (CC BY 4.0). This material is distributed under the Creative Commons Attribution 4.0 International licence, http://creativecommons.org/licenses/by/4.0/, as recorded in the arXiv OAI licence metadata for this exact version and shown on the abstract page https://arxiv.org/abs/2607.00061v1. Full rights notice: licenses/arxiv-2607.00061-CC-BY-4.0.txt.\par\smallskip
\noindent\textit{Editorial scope.} Counted unit: the original \texttt{theorem} environment \texttt{thm:6.1} at source lines 494--504 together with its complete proof at lines 506--521; the delivered document keeps source lines 184--641 so that every referenced label and every citation resolves inside it. Native editable source is the paper's own concatenated TeX; the AMS wrapper is editorial formatting only. No publisher class, upstream script or unrelated artwork is redistributed.
\begin{equation} \label{2-1}
\begin{cases}
\dfrac{dT}{dt} = \lambda_T - \mu_T T - \chi T V, \\[6pt]
\dfrac{dI}{dt} = \chi T V - \mu_I I - \alpha I Z_a, \\[6pt]
\dfrac{dV}{dt} = \varepsilon_V \mu_I I - \mu_V V, \\[6pt]
\dfrac{dZ}{dt} = \lambda_Z - \mu_Z Z - \beta Z I, \\[6pt]
\dfrac{dZ_a}{dt} = \beta Z I - \mu_{Z_a} Z_a .
\end{cases}
\end{equation}
The model parameters are as follows: the natural mortality rate of healthy CD4\(^+\) T-cells, denoted \(\mu_T\), the natural death rate of infected CD4\(^+\) T-cells, also denoted \(\mu_T\), the infection transmission rate of CD4\(^+\) T-cells upon contact with free virions, \(\lambda\). The clearance rate of free infectious virions, \(\mu_V\). The constant recruitment rate of CD8\(^+\) T-cells into the system, \(\lambda_Z\), the removal rate of infected CD4\(^+\) T-cells mediated by activated CD8\(^+\) T-cells, \(\alpha\),and the natural mortality rate of CD8\(^+\) T-cells, denoted \(\mu_Z\). The natural decay rate of activated CD8\(^+\) (immune response) cells, also indicated by \(\mu_Z\), and the activation rate of CD8\(^+\) T-cells stimulated by the presence of infected CD4\(^+\) T-cells, \(\beta\). The average production rate of HIV virions by each infected CD4\(^+\) T-cell, \(c_V\).
%%%======================================================================================================================
\section{Preliminaries}\label{sec3}
\begin{definition}\label{def2-1}
To define the fractal-fractional derivative of order $\nu_1$ in the Riemann-Liouville terms of functions $\beta(\tilde{t})$ that are not essential differentiable. Where $(\gamma)$ fractional dimension.
Then for $0 \leq
\nu_1,~ \eta \leq 1$, .
\begin{itemize}
    \item the power law kernel is defined as:
\begin{equation}\label{eq2-1}
    _0^{FFP}D_t^{\nu_1, \gamma}  \Phi (\tilde{t}) = \frac{1}{\Gamma (n-\nu_1)} \frac{d^n}{dt^n} \int_0^{\tilde{t}} (t-\zeta_1)^{n-\nu_1-1} \Phi (\zeta_1) d\zeta_1
\end{equation}
where $n-1 < \nu_1,~ \gamma <n \in N$. And 
\begin{equation}\label{eq2-2}
\frac{D \Phi(\zeta_1)}{D \zeta_1^{\eta}} = \lim_{t \rightarrow \zeta_1} \frac{\psi(\tilde{t}) - \Phi(\zeta_1)}{ t^{\eta} - \zeta_1^{\eta}}   
\end{equation}
\item the exponential decay kernel is obtained as:
\begin{equation}\label{eq2-3}
    _0^{FFE}D_t^{\nu_1, \gamma}  \Phi (\tilde{t}) = \frac{M (\nu_1)}{\Gamma (n-\nu_1)} \frac{d^n}{dt^n} \int_0^{\tilde{t}} exp [- \frac{\nu_1}{1-\nu_1}(t-\zeta_1)] \Phi (\zeta_1) d\zeta_1    
\end{equation}

\item with the application of the Mittag-Leffler kernel (MLK) as fellow:
\begin{equation}\label{eq2-4}
    _0^{FFM}D_t^{\nu_1, \gamma}  \Phi (\tilde{t}) = \frac{AB (\nu_1)}{1-\nu_1 } \frac{d^n}{dt^n} \int_0^{\tilde{t}} \Phi (\zeta_1) E_{\nu_1} [- \frac{\nu_1}{1-\nu_1}(t-\zeta_1)^{\nu_1}]  d\zeta_1    
\end{equation}
\end{itemize}
\end{definition}
%--------------------------------------------------------------------------
\begin{definition}\label{def2-2}
For a continuous function $\Phi(\tilde{t})$ on $(a,b)$, then the fractal-fractional
integral of $\phi(\tilde{t})$.
\begin{itemize}
    \item power law kernel is epressed as:
\begin{equation}\label{eq2-5}
    _0^{FFP}I^{\nu_1, \gamma}  \Phi (\tilde{t}) = \frac{1}{\Gamma ( \nu_1)} \int_0^{\tilde{t}} (t-\zeta_1)^{\nu_1-1} \zeta_1^{1-\eta} \Phi (\zeta_1) d\zeta_1
\end{equation}
\item the exponential decay kernel is represented as:
\begin{equation}\label{eq2-6}
    _0^{FFE}I^{\nu_1, \gamma}  \Phi (\tilde{t}) = \frac{\eta (1-\nu_1) t^{\eta -1} \Phi(\tilde{t})}{M(\nu_1)} +  \frac{\nu_1 \eta}{M (\nu_1)}  \int_0^{\tilde{t}} \zeta_1^{\nu_1-1} \Phi (\zeta_1) d\zeta_1    
\end{equation}
\item similarly MLK is:
\begin{equation}\label{eq2-7}
    _0^{FFM}I^{\nu_1, \gamma}  \Phi (\tilde{t}) = \frac{\eta (a-\nu_1) t^{\eta -1} \beta(\tilde{t})}{AB (\eta_1)} + \frac{ \nu_1 \eta }{AB(\nu_1) \Gamma(\nu_1) } \int_0^{\tilde{t}} (t- \zeta_1)^{\nu_1 -1 } \zeta_1^{\eta-1} \Phi (\zeta_1)   d\zeta_1    
\end{equation}
\end{itemize}
\end{definition}
%========================================================================================
\section{Fractal--Fractional Model Formulation}\label{sec4}
We extend the classical HIV (\ref{2-1})infection model is transformed using MLK to a fractal--fractional framework. This formulation captures memory effects and heterogeneity in the immune response dynamics.
%\subsection{Fractal--Fractional Derivative with Mittag--Leffler Kernel}
Let $0<\alpha<1$ denote the fractional order and $0<\gamma\leq 1$ the fractal dimension. The fractal--fractional derivative in the Caputo sense with MKL (Atangana--Baleanu type) of a sufficiently smooth function $f(\tilde{t})$ is defined by
\begin{equation}
{}^{FFM}D_{t}^{\alpha,\gamma} f(\tilde{t})
=
\frac{B(\alpha)}{1-\alpha}
\int_{0}^{t}
E_{\alpha}\!\left(-\frac{\alpha}{1-\alpha}(t-\tau)^{\alpha}\right)
\frac{d}{d\tau}\big(f(\tau)\big)\,
\tau^{1-\gamma}\, d\tau,
\end{equation}
where $E_{\alpha}(\cdot)$ denotes the Mittag--Leffler function and $B(\alpha)$ is a normalization constant satisfying $B(0)=B(1)=1$.
Using the above operator, the fractal--fractional HIV infection model (\ref{eq2-1}) is formulated as
\begin{equation}
\label{eq:FF_HIV_model}
\begin{cases}
{}^{FFM}D_{t}^{\alpha,\gamma} T(\tilde{t})
= \lambda_T - (\mu_T + \chi V(\tilde{t}))T(\tilde{t}), \\[6pt]

{}^{FFM}D_{t}^{\alpha,\gamma} I(\tilde{t})
= \chi T(\tilde{t})V(\tilde{t}) - (\mu_I  + \alpha Z_{a}(\tilde{t}))I(\tilde{t}), \\[6pt]

{}^{FFM}D_{t}^{\alpha,\gamma} V(\tilde{t})
= \varepsilon_V \mu_I I(\tilde{t}) - \mu_V V(\tilde{t}), \\[6pt]

{}^{FFM}D_{t}^{\alpha,\gamma} Z(\tilde{t})
= \lambda_Z - (\mu_Z  + \beta  I(\tilde{t}))Z(\tilde{t}), \\[6pt]

{}^{FFM}D_{t}^{\alpha,\gamma} Z_{a}(\tilde{t})
= \beta Z(\tilde{t}) I(\tilde{t}) - \mu_{Z_a} Z_{a}(\tilde{t}).
\end{cases}
\end{equation}
The model \eqref{eq:FF_HIV_model} is supplemented with the following nonnegative history:
\begin{equation}\label{eq:initial}
(T(0),I(0),V(0),Z(0), Z_a(0) )=(T_0,I_0,V_0,Z_0,Z_{a0} )
\end{equation}
where $T_0$, $I_0$, $V_0$, $Z_0$, and $Z_{a0}$ are given constants.
%===========================================================================================
\section{Proposed Model's Existence and Individuality}
Now, we develop, and prove the basic mathematical properties of our proposed model \eqref{eq:FF_HIV_model}. By changing our proposed model (\ref{eq:FF_HIV_model}) into its integral form, and utilize the fixed-point theorem.

\subsection{Existence of the Solution}
We first reformulate the system \eqref{eq:FF_HIV_model} into an analogous integral form with the application of the fractal-fractional integral operator, let us assume $K(t)= \frac{(\alpha-1)\gamma}t^{\gamma-1}{\mathcal{AB(\alpha)}\Gamma(\alpha)}$, and $L = \frac{\alpha\gamma}{\mathcal{AB(\alpha)}\Gamma(\alpha)}$ \\
\begin{equation}
	\label{eq:integral_system}
	\begin{cases}
		T(\tilde{t}) = T(0) + K(\tilde{t})H_1(\tilde{t}, T) + L \int_0^{\tilde{t}} (\tilde{t}-\zeta_1)^{\alpha-1} \zeta_1^{\gamma-1} H_1(\zeta_1, T) d\zeta_1, \\
		I(\tilde{t}) = I(0) + K(\tilde{t}) H_2(\tilde{t}, I) + L \int_0^{\tilde{t}} (\tilde{t}-\zeta_1)^{\alpha-1} \zeta_1^{\gamma-1} H_2(\zeta_1, I) d\zeta_1, \\
		V(\tilde{t}) = V(0) + K(\tilde{t}) H_3(\tilde{t}, V) + L \int_0^{\tilde{t}} (\tilde{t}-\zeta_1)^{\alpha-1} \zeta_1^{\gamma-1} H_3(\zeta_1, V) d\zeta_1, \\
		Z(\tilde{t}) = Z(0) + K(\tilde{t}) H_4(\tilde{t}, Z) + L \int_0^{\tilde{t}} (\tilde{t}-\zeta_1)^{\alpha-1} \zeta_1^{\gamma-1} H_4(\zeta_1, Z) d\zeta_1, \\
		Z_{a}(\tilde{t}) = Z_a(0) + K(\tilde{t}) H_5(\tilde{t}, Z_a) + L \int_0^{\tilde{t}} (\tilde{t}-\zeta_1)^{\alpha-1} \zeta_1^{\gamma-1} H_5(\zeta_1, Z_a) d\zeta_1.
	\end{cases}
\end{equation}
where the kernels $H_j$ represent the right-hand side of the model equations. We define an iterative sequence $T_n(\tilde{t}), I_n(\tilde{t}), V_n(\tilde{t}), Z_n(\tilde{t}), Z_{an}(\tilde{t})$ to show convergence.
\subsection{Definition of the Kernels}
The kernels $H_i(\tilde{t}, \beta(\tilde{t}))$ for $i \in \{1, 2, 3, 4, 5\}$, show the right hand side of the proposed HIV model \eqref{eq:FF_HIV_model}, and $ X = (t, T, I, V, Z, Z_a)$.
In the subsequent existence, and uniqueness analysis, we assume that the state variables are bounded within a Banach space $C[0, b]$. Therefore, there exist positive constants $A_i$ such that:
\begin{theorem}
Assume that the continuous functions \(T, I, V, Z, Z_a \in L[0,1]\) satisfy 
\(\|T\| \leq \phi_1,\; \|I\| \leq \phi_2,\; \|V\| \leq \phi_3,\; \|Z\| \leq \phi_4,\; \|Z_a\| \leq \phi_5\) 
with each \(\phi_j < 1\). Then each kernel \(H_k\;(m=1,\dots,5)\) is Lipschitz and a contraction.
\end{theorem}
\begin{proof}
For two functions $T(\tilde{\tilde{t}})$ and $\hat{T}(\tilde{t})$, we have:
	\begin{align*}
		\|H_1(\tilde{t}, T) - H_1(\tilde{t}, \hat{T})\| &= \|(\lambda_T - \mu_T T - \chi TV) - (\lambda_T - \mu_T \hat{T} - \chi \hat{T}V)\| \\
		&= \|-\mu_T(T - \hat{T}) - \chi V(T - \hat{T})\| \\
		&\leq [\mu_T + \chi \|V\|] \|T - \hat{T}\| \\
		&\leq [\mu_T + \chi A_3] \|T - \hat{T}\| \\
		&\leq \phi_1 \|T - \hat{T}\|,
	\end{align*}
	where $\phi_1 = \mu_T + \chi A_3 < 1$. Thus, $H_1$ fulfills the Lipschitz condition.
	
	Next, for $H_2(\tilde{t}, I)$ with $I(\tilde{t})$ and $\hat{I}(\tilde{t})$:
	\begin{align*}
		\|H_2(\tilde{t}, I) - H_2(\tilde{t}, \hat{I})\| &= \|(\chi TV - \mu_I I - \alpha I Z_a) - (\chi TV - \mu_I \hat{I} - \alpha \hat{I} Z_a)\| \\
		&\leq [\mu_I + \alpha \|Z_a\|] \|I - \hat{I}\| \\
		&\leq [\mu_I + \alpha A_5] \|I - \hat{I}\| \\
		&\leq \phi_2 \|I - \hat{I}\|,
	\end{align*}
	where $\phi_2 = \mu_I + \alpha A_5 < 1$.
	
	For $H_3(\tilde{t}, V)$ with $V(\tilde{t})$ and $\hat{V}(\tilde{t})$:
	\begin{align*}
		\|H_3(\tilde{t}, V) - H_3(\tilde{t}, \hat{V})\| &= \|(\varepsilon_V \mu_I I - \mu_V V) - (\varepsilon_V \mu_I I - \mu_V \hat{V})\| \\
		&\leq \mu_V \|V - \hat{V}\| \\
		&\leq \phi_3 \|V - \hat{V}\|,
	\end{align*}
	where $\phi_3 = \mu_V < 1$.
	
	For $H_4(\tilde{t}, Z)$ with $Z(\tilde{t})$ and $\hat{Z}(\tilde{t})$:
	\begin{align*}
		\|H_4(\tilde{t}, Z) - H_4(\tilde{t}, \hat{Z})\| &= \|(\lambda_Z - \mu_Z Z - \beta Z I) - (\lambda_Z - \mu_Z \hat{Z} - \beta \hat{Z} I)\| \\
		&\leq [\mu_Z + \beta \|I\|] \|Z - \hat{Z}\| \\
		&\leq [\mu_Z + \beta A_2] \|Z - \hat{Z}\| \\
		&\leq \phi_4 \|Z - \hat{Z}\|,
	\end{align*}
	where $\phi_4 = \mu_Z + \beta A_2 < 1$.
	
	Finally, for $H_5(\tilde{t}, Z_a)$ with $Z_{a}(\tilde{t})$ and $\hat{Z}_a(\tilde{t})$:
	\begin{align*}
		\|H_5(\tilde{t}, Z_a) - H_5(\tilde{t}, \hat{Z}_a)\| &= \|(\beta Z I - \mu_{Z_a} Z_a) - (\beta Z I - \mu_{Z_a} \hat{Z}_a)\| \\
		&\leq \mu_{Z_a} \|Z_a - \hat{Z}_a\| \\
		&\leq \phi_5 \|Z_a - \hat{Z}_a\|,
	\end{align*}
	where $\phi_5 = \mu_{Z_a} < 1$.
    When $\phi_j < 1$ then all the kernels verified the desire conditions.  the proof is completed.
\end{proof}
% \subsection{Analysis of existence and uniqueness of the solutions}
% In this subsection, we demonstrate the theorem that establishes the existence and distinctiveness of solutions for the proposed HIV model \eqref{eq:FF_HIV_model}. We verify the existence of the fractal-fractional model using the fixed-point approach and Lipschitz condition.



\subsection{Iterative Scheme for Qualitative Analysis}
Using of the \eqref{eq:integral_system} to determine the existence of solutions, we define the following iterative convergent sequence:
\begin{equation}
	\begin{cases}
		T_n(\tilde{t}) = K(\tilde{t}) H_1(\tilde{t}, T_{n-1}) + L \int_0^{\tilde{t}} (\tilde{t}-\zeta_1)^{\alpha-1} \zeta_1^{\gamma-1} H_1(\zeta_1, T_{n-1}) d\zeta_1, \\[10pt]
		I_n(\tilde{t}) = K(\tilde{t}) H_2(\tilde{t}, I_{n-1}) + L \int_0^{\tilde{t}} (\tilde{t}-\zeta_1)^{\alpha-1} \zeta_1^{\gamma-1} H_2(\zeta_1, I_{n-1}) d\zeta_1, \\[10pt]
		V_n(\tilde{t}) = K(\tilde{t}) H_3(\tilde{t}, V_{n-1}) + L \int_0^{\tilde{t}} (\tilde{t}-\zeta_1)^{\alpha-1} \zeta_1^{\gamma-1} H_3(\zeta_1, V_{n-1}) d\zeta_1, \\[10pt]
		Z_n(\tilde{t}) = K(\tilde{t}) H_4(\tilde{t}, Z_{n-1}) + L \int_0^{\tilde{t}} (\tilde{t}-\zeta_1)^{\alpha-1} \zeta_1^{\gamma-1} H_4(\zeta_1, Z_{n-1}) d\zeta_1, \\[10pt]
		Z_{an}(\tilde{t}) = K(\tilde{t}) H_5(\tilde{t}, Z_{an-1}) + L \int_0^{\tilde{t}} (\tilde{t}-\zeta_1)^{\alpha-1} \zeta_1^{\gamma-1} H_5(\zeta_1, Z_{an-1}) d\zeta_1.
	\end{cases}
\end{equation}
Next, we consider the differences between successive terms of the iterative sequence as follows:
\begin{align*}
	\Delta T_{n}(\tilde{t}) &= T_{n}(\tilde{t}) - T_{n-1}(\tilde{t}) \\
	&= K(\tilde{t}) \left( H_1(\tilde{t}, T_{n-1}) - H_1(\tilde{t}, T_{n-2}) \right) \\
	&\quad + L \int_0^{\tilde{t}} (\tilde{t}-\zeta_1)^{\alpha-1} \zeta_1^{\gamma-1} \left( H_1(\zeta_1, T_{n-1}) - H_1(\zeta_1, T_{n-2}) \right) d\zeta_1, \\[10pt]
	\Delta I_{n}(\tilde{t}) &= I_{n}(\tilde{t}) - I_{n-1}(\tilde{t}) \\
	&= K(\tilde{t}) \left( H_2(t, I_{n-1}) - H_2(\tilde{t}, I_{n-2}) \right) \\
	&\quad + L \int_0^{\tilde{t}} (\tilde{t}-\zeta_1)^{\alpha-1} \zeta_1^{\gamma-1} \left( H_2(\zeta_1, I_{n-1}) - H_2(\zeta_1, I_{n-2}) \right) d\zeta_1, \\[10pt]
	\Delta V_{n}(\tilde{t}) &= V_{n}(\tilde{t}) - V_{n-1}(\tilde{t}) \\
	&= K(\tilde{t}) \left( H_3(\tilde{t}, V_{n-1}) - H_3(\tilde{t}, V_{n-2}) \right) \\
	&\quad + L \int_0^{\tilde{t}} (\tilde{t}-\zeta_1)^{\alpha-1} \zeta_1^{\gamma-1} \left( H_3(\zeta_1, V_{n-1}) - H_3(\zeta_1, V_{n-2}) \right) d\zeta_1, \\[10pt]
\end{align*}
   \begin{align*} 
	\Delta Z_{n}(\tilde{t}) &= Z_{n}(\tilde{t}) - Z_{n-1}(\tilde{t}) \\
	&= K(\tilde{t}) \left( H_4(\tilde{t}, Z_{n-1}) - H_4(\tilde{t}, Z_{n-2}) \right) \\
	&\quad + L \int_0^{\tilde{t}} (\tilde{t}-\zeta_1)^{\alpha-1} \zeta_1^{\gamma-1} \left( H_4(\zeta_1, Z_{n-1}) - H_4(\zeta_1, Z_{n-2}) \right) d\zeta_1, \\[10pt]
	\Delta Z_{an}(\tilde{t}) &= Z_{an}(\tilde{t}) - Z_{an-1}(\tilde{t}) \\
	&= K(\tilde{t}) \left( H_5(\tilde{t}, Z_{an-1}) - H_5(\tilde{t}, Z_{an-2}) \right) \\
	&\quad + L \int_0^{\tilde{t}} (\tilde{t}-\zeta_1)^{\alpha-1} \zeta_1^{\gamma-1} \left( H_5(\zeta_1, Z_{an-1}) - H_5(\zeta_1, Z_{an-2}) \right) d\zeta_1.
\end{align*}
Taking the norms of the aforementioned system on both sides, we obtained:
\begin{align*}
	\|\Delta T_{n+1}(\tilde{t})\| &= K(\tilde{t}) \|H_1(\tilde{t}, T_n) - H_1(\tilde{t}, T_{n-1})\| \\
	&\quad + L \int_0^{\tilde{t}} (\tilde{t}-\zeta_1)^{\alpha-1} \zeta_1^{\gamma-1} \|H_1(\zeta_1, T_n) - H_1(\zeta_1, T_{n-1})\| d\zeta_1, \\[10pt]
	\|\Delta I_{n+1}(\tilde{t})\| &= K(\tilde{t}) \|H_2(\tilde{t}, I_n) - H_2(\tilde{t}, I_{n-1})\| \\
	&\quad + L \int_0^{\tilde{t}} (\tilde{t}-\zeta_1)^{\alpha-1} \zeta_1^{\gamma-1} \|H_2(\zeta_1, I_n) - H_2(\zeta_1, I_{n-1})\| d\zeta_1, \\[10pt]
	\|\Delta V_{n+1}(\tilde{t})\| &= K(\tilde{t}) \|H_3(\tilde{t}, V_n) - H_3(\tilde{t}, V_{n-1})\| \\
	&\quad + L \int_0^{\tilde{t}} (\tilde{t}-\zeta_1)^{\alpha-1} \zeta_1^{\gamma-1} \|H_3(\zeta_1, V_n) - H_3(\zeta_1, V_{n-1})\| d\zeta_1, \\[10pt]
	\|\Delta Z_{n+1}(\tilde{t})\| &= K(\tilde{t}) \|H_4(\tilde{t}, Z_n) - H_4(\tilde{t}, Z_{n-1})\| \\
	&\quad + L \int_0^{\tilde{t}} (\tilde{t}-\zeta_1)^{\alpha-1} \zeta_1^{\gamma-1} \|H_4(\zeta_1, Z_n) - H_4(\zeta_1, Z_{n-1})\| d\zeta_1, \\[10pt]
	\|\Delta Z_{an+1}(\tilde{t})\| &= K(\tilde{t}) \|H_5(\tilde{t}, Z_{an}) - H_5(\tilde{t}, Z_{an-1})\| \\
	&\quad + L \int_0^{\tilde{t}} (\tilde{t}-\zeta_1)^{\alpha-1} \zeta_1^{\gamma-1} \|H_5(\zeta_1, Z_{an}) - H_5(\zeta_1, Z_{an-1})\| d\zeta_1.
\end{align*}

By applying the Lipschitz conditions $\|H_j(\Psi_n) - H_j(\Psi_{n-1})\| \leq \phi_j \|\Delta \Psi_n\|$, we obtain the generalized bounding system by taking $C(t) =\left[ K(\tilde{t}) + \frac{\alpha\gamma \Gamma(\gamma) t^{\alpha+\gamma-1}}{\mathcal{AB(\alpha)}\Gamma(\alpha+\gamma)} \right] $.
\begin{align*}
	\|\Delta T_{n+1}(\tilde{t})\| &\leq C(\tilde{t}) \phi_1 \|\Delta T_n(\tilde{t})\|, \\[10pt]
	\|\Delta I_{n+1}(\tilde{t})\| &\leq C(\tilde{t}) \phi_2 \|\Delta I_n(\tilde{t})\|, \\[10pt]
	\|\Delta V_{n+1}(\tilde{t})\| &\leq C(\tilde{t}) \phi_3 \|\Delta V_n(\tilde{t})\|, \\[10pt]
	\|\Delta Z_{n+1}(\tilde{t})\| &\leq C(\tilde{t}) \phi_4 \|\Delta Z_n(\tilde{t})\|, \\[10pt]
	\|\Delta Z_{an+1}(\tilde{t})\| &\leq C(\tilde{t}) \phi_5 \|\Delta Z_{an}(\tilde{t})\|.
\end{align*}
\begin{theorem}
	The solution of our proposed model \eqref{eq:FF_HIV_model} is exist if:
	\begin{equation}
		\Omega = \max[\phi_1, \phi_2, \phi_3, \phi_4, \phi_5] < 1.
	\end{equation}
\end{theorem}

\begin{proof}
	We define the difference functions between the iterative sequence and the actual solution as follows:
	\begin{equation*}
		\begin{cases}
			\mathcal{Z}_{1n}(\tilde{t}) = T_{n+1}(\tilde{t}) - T(\tilde{t}), \\
			\mathcal{Z}_{2n}(\tilde{t}) = I_{n+1}(\tilde{t}) - I(\tilde{t}), \\
			\mathcal{Z}_{3n}(\tilde{t}) = V_{n+1}(\tilde{t}) - V(\tilde{t}), \\
			\mathcal{Z}_{4n}(\tilde{t}) = Z_{n+1}(\tilde{t}) - Z(\tilde{t}), \\
			\mathcal{Z}_{5n}(\tilde{t}) = Z_{a,n+1}(\tilde{t}) - Z_{a}(\tilde{t}).
		\end{cases}
	\end{equation*}
	
	Applying the norm of the first equation $\mathcal{Z}_{1n}(\tilde{t})$ and applying the Lipschitz criterion $\phi_1$, we get:
	\begin{align*}
		\|\mathcal{Z}_{1n}(\tilde{t})\| &= \|T_{n+1}(\tilde{t}) - T(\tilde{t})\| \\
		&\leq C(\tilde{t}) \phi_1 \|T_n - T\| \\
		&\leq C(\tilde{t})^n \lambda^n \|T_n - T\|,
	\end{align*}
	where $\lambda < 1$. As $n \to \infty$, we observe that $\|\mathcal{Z}_{1n}(\tilde{t})\| \to 0$. Similarly, for the remaining compartments:
	\begin{align*}
		\|\mathcal{Z}_{2n}\| &\leq C(\tilde{t})^n \lambda^n \|I_n - I\| \to 0, \\
		\|\mathcal{Z}_{3n}\| &\leq C(\tilde{t})^n \lambda^n \|V_n - V\| \to 0, \\
		\|\mathcal{Z}_{4n}\| &\leq C(\tilde{t})^n \lambda^n \|Z_n - Z\| \to 0, \\
		\|\mathcal{Z}_{5n}\| &\leq C(\tilde{t})^n \lambda^n \|Z_{an} - Z_a\| \to 0.
	\end{align*}
	
	Which show that $\mathcal{Z}_{jn}(\tilde{t}) \to 0$ as $n \to \infty$ for $j \in \{1, 2, 3, 4, 5\}$, which confirms that a solution exists for the HIV model.
\end{proof}
\begin{theorem}
	If $C(\tilde{t}) \phi_j < 1, \quad j \in \{1, 2, 3, 4, 5\}$, then Then our proposed model \eqref{eq:FF_HIV_model} has a unique solution.
\end{theorem}

\begin{proof}
 Let $\{T_1, I_1, V_1, Z_1, Z_a\}$ and $\{\tilde{T_1}, \tilde{I_1}, \tilde{V_1}, \tilde{Z_1}, \tilde{Z}_a\}$ be two distinct solutions of the system. 
	
	Taking the difference of $T_1(\tilde{t})$, and $\tilde{T_1}(\tilde{t})$ then applying the norm, we obtain:
	\begin{align*}
		\|T_1(\tilde{t}) - \tilde{T_1}(\tilde{t})\| &= \Bigg\| K(\tilde{t}) \left( H_1(\tilde{t}, T) - H_1(\tilde{t}, \tilde{T}) \right) \\
		&\quad + L \int_0^{\tilde{t}} (\tilde{t}-\zeta_1)^{\alpha-1} \zeta_1^{\gamma-1} \left( H_1(\zeta_1, T) - H_1(\zeta_1, \tilde{T}) \right) d\zeta_1 \Bigg\|.
	\end{align*}
	
	Applying the triangle inequality and the Lipschitz criterion $\phi_1$:
	\begin{align*}
		\|T_1(\tilde{t}) - \tilde{T_1}(\tilde{t})\| &\leq K(\tilde{t}) \phi_1 \|T - \tilde{T}\| \\
		&\quad + L \phi_1 \|T_1 - \tilde{T_1}\| \int_0^{\tilde{t}} (\tilde{t}-\zeta_1)^{\alpha-1} \zeta_1^{\gamma-1} d\zeta_1.
	\end{align*}
	
	Evaluating the integral $\int_0^{\tilde{t}} (\tilde{t}-\zeta_1)^{\alpha-1} \zeta_1^{\gamma-1} d\zeta_1 = \frac{\Gamma(\alpha)\Gamma(\gamma)}{\Gamma(\alpha+\gamma)} t^{\alpha+\gamma-1}$, we obtain:
	\begin{equation*}
		\|T_1(\tilde{t}) - \tilde{T_1}(\tilde{t})\| \leq C(\tilde{t}) \phi_1 \|T_1 - \tilde{T_1}\|.
	\end{equation*}
	
	Rearranging the terms leads to:
	\begin{equation*}
		\|T_1 - \tilde{T_1}\| - C(\tilde{t}) \phi_1 \|T_1 - \tilde{T_1}\| \leq 0,
	\end{equation*}
	\begin{equation}
		( 1 - C(\tilde{t}) )\|T_1 - \tilde{T_1}\| \leq 0.
	\end{equation}
	
	The inequality mentioned above is accurate if and only if:
	\begin{equation*}
		\|T_1 - \tilde{T_1}\| = 0 \implies T_1 = \tilde{T_1}.
	\end{equation*}
	
	Similarly, by repeating the same steps for the other compartments, we obtain $I_1 = \tilde{I_1},V_1 = \tilde{V_1},Z_1 = \tilde{Z_1},$ and $Z_a = \tilde{Z}_a.$
	Therefore, the system \eqref{eq:FF_HIV_model} has a unique solution.
\end{proof}
\subsection{Hyers-Ulam Stability Analysis}
In this subsection, we investigate the Hyers-Ulam (HU) stability of the proposed fractal-fractional HIV model \eqref{eq:FF_HIV_model}. We begin by defining the stability criteria for the system.
\subsection{Basic Definitions}
Let $\epsilon > 0$ be a small positive constant, and let the following inequality hold for the fractal-fractional operator:
\begin{equation}
	\left\| {}^{FFM}D_{\tilde{t}}^{\alpha,\gamma} \beta(\tilde{t}) - H_j(\tilde{t}, \beta(\tilde{t})) \right\| \leq \epsilon, \quad j \in \{1, \dots, 5\}.
\end{equation}
A solution $\beta(\tilde{t})$ of the model is said to be Hyers-Ulam stable if there exists a constant $C > 0$ and an exact solution $\beta^*(\tilde{t})$ such that:
\begin{equation}
	\|\beta(\tilde{t}) - \beta^*(\tilde{t})\| \leq C \epsilon.
\end{equation}

\begin{theorem}\label{thm:6.1}
Let $H_j$ be the kernels of the system \eqref{eq:FF_HIV_model}. If $\beta(\tilde{t})$ is an approximate solution satisfying the inequality:
\begin{equation}
    \left\| {}^{FFM}D_{\tilde{t}}^{\alpha,\gamma} \beta
    (\tilde{t}) - H_j(\tilde{t}, \beta(\tilde{t})) \right\| \leq \epsilon, \quad \epsilon > 0,
\end{equation}
then $\beta(\tilde{t})$ satisfies the following integral inequality:
\begin{equation}
    \left\| \beta(\tilde{t}) - \left( \beta(0) + \mathcal{I}_{FF}^{\alpha, \gamma} H_j(\tilde{t}, \beta(\tilde{t})) \right) \right\| \leq C(\tilde{t}) \epsilon = \Lambda \epsilon.
\end{equation}
\end{theorem}

\begin{proof}
By integrating the perturbed differential inequality, we obtain:
\begin{align*}
    \beta(\tilde{t}) &\leq \beta(0) + K(\tilde{t}) (H_j(t, \Psi(\tilde{t})) + \epsilon) \\
    &\quad + L \int_0^{\tilde{t}} (\tilde{t}-\zeta_1)^{\alpha-1} \zeta_1^{\gamma-1} (H_j(\zeta_1, \beta(\zeta_1)) + \epsilon) d\zeta_1.
\end{align*}
Applying the norm and subtracting the exact integral form $\beta(0) + \mathcal{I}_{FF}^{\alpha, \gamma} H_j(\tilde{t}, \beta(\tilde{t}))$ from both sides:
\begin{align*}
    \left\| \Psi(\tilde{t}) - \left( \beta(0) + \mathcal{I}_{FF}^{\alpha, \gamma} H_j(\tilde{t}, \Psi(\tilde{t})) \right) \right\| &\leq \epsilon \left[ K(\tilde{t}) + L \int_0^{\tilde{t}} (\tilde{t}-\zeta_1)^{\alpha-1} \zeta_1^{\gamma-1} d\zeta_1 \right] \\
    &\leq \epsilon C(\tilde{t}).
\end{align*}
By defining the term as $\Lambda = C(\tilde{t})$, we arrive at:
\begin{equation*}
    \left\| \beta(\tilde{t}) - \left( \beta(0) + \mathcal{I}_{FF}^{\alpha, \gamma} H_j(\tilde{t}, \beta(\tilde{t})) \right) \right\| \leq \Lambda \epsilon.
\end{equation*}
\end{proof}

\begin{theorem}
	The fractal-fractional HIV model \eqref{eq:FF_HIV_model} is Hyers-Ulam stable if the condition $\phi_j \Lambda < 1$ holds for all $j \in \{1, 2, 3, 4, 5\}$, where $\phi_j$ are the Lipschitz constants and $\Lambda$ is the integral bound defined in Theorem \ref{thm:6.1}.
\end{theorem}

\begin{proof}
	Let $\beta(\tilde{t}) = \{T-1, I_1, V_1, Z_1, Z_a\}$ be the approximate solution of the system satisfying the perturbed inequality, and let $\beta^*(\tilde{t}) = \{T_1^*, I_1^*, V_1^*, Z_1^*, Z_a^*\}$ be the unique exact solution of the system \eqref{eq:FF_HIV_model}.
	
	Starting with the first compartment, $T(\tilde{t})$, we consider the difference:
	\begin{align*}
		\|T_1(\tilde{t}) - T_1^*(\tilde{t})\| &= \left\| T_1(\tilde{t}) - \left( T_(0) + \mathcal{I}_{FF}^{\alpha, \gamma} H_1(\tilde{t}, T_1^*) \right) \right\| \\
		&\leq \left\| T_1(\tilde{t}) - \left( T_1(0) + \mathcal{I}_{FF}^{\alpha, \gamma} H_1(\tilde{t}, T_1) \right) \right\| + \left\| \mathcal{I}_{FF}^{\alpha, \gamma} H_1(\tilde{t}, T_1) - \mathcal{I}_{FF}^{\alpha, \gamma} H_1(\tilde{t}, T_1^*) \right\|.
	\end{align*}
	
	Using the result from Theorem \ref{thm:6.1} for the first term and the Lipschitz condition for the second term:
	\begin{equation*}
		\|T_1(\tilde{t}) - T_1^*(\tilde{t})\| \leq \Lambda \epsilon + C(\tilde{t}) \phi_1 \|T_1(\tilde{t}) - T_1^*(\tilde{t})\|.
	\end{equation*}
	
	Applying the notation $\Lambda = K(\tilde{t}) + \frac{\alpha\gamma \Gamma(\gamma) t^{\alpha+\gamma-1}}{\mathcal{AB(\alpha)}\Gamma(\alpha+\gamma)}$, the inequality becomes:
	\begin{equation*}
		\|T_1(\tilde{t}) - T_1^*(\tilde{t})\| \leq \Lambda \epsilon + \phi_1 \Lambda \|T_1(\tilde{t}) - T_1^*(\tilde{t})\|.
	\end{equation*}
	
	Rearranging the terms to isolate the norm of the difference:
	\begin{equation*}
		\|T_1(\tilde{t}) - T_1^*(\tilde{t})\| (1 - \phi_1 \Lambda) \leq \Lambda \epsilon,
	\end{equation*}
	which leads to:
	\begin{equation}
		\|T_1(\tilde{t}) - T_1^*(\tilde{t})\| \leq \frac{\Lambda}{1 - \phi_1 \Lambda} \epsilon = C_1 \epsilon.
	\end{equation}
	
	By repeating the same logical steps for the remaining compartments $I_1, V_1, Z_1,$ and $Z_a$, we obtain:
	\begin{align*}
		\|I_1(\tilde{t}) - I_1^*(\tilde{t})\| &\leq C_2 \epsilon, \\
		\|V_1(\tilde{t}) - V_1^*(\tilde{t})\| &\leq C_3 \epsilon, \\
		\|Z_1(\tilde{t}) - Z_1^*(\tilde{t})\| &\leq C_4 \epsilon, \\
		\|Z_{a}(\tilde{t}) - Z_a^*(\tilde{t})\| &\leq C_5 \epsilon,
	\end{align*}
	where $C_j = \frac{\Lambda}{1 - \phi_j \Lambda} > 0$. Since all differences are bounded by a constant multiple of $\epsilon$, the proposed fractal-fractional HIV model is Hyers-Ulam stable.
\end{proof}
\section{Numerical Methodology}\label{sec7}
In order to numerically solve the proposed fractal-fractional HIV model \eqref{eq:FF_HIV_model} under the Mittag-Leffler kernel framework, we describe a numerical technique based on a Newton polynomial method \cite{7}. This approach captures non-local historical memory profiles and fractal characteristics simultaneously. For convenience and clarity of the formulation, we rewrite the system \eqref{eq:FF_HIV_model} as follows:
\begin{equation}
\label{eq:functional_form}
\begin{cases}
{}^{FFM}_{0}D_{t}^{\alpha,\gamma} T(\tilde{t}) = \mathcal{F}_1X, \\[6pt]
{}^{FFM}_{0}D_{t}^{\alpha,\gamma} I(\tilde{t}) = \mathcal{F}_2X, \\[6pt]
{}^{FFM}_{0}D_{t}^{\alpha,\gamma} V(\tilde{t}) = \mathcal{F}_3X, \\[6pt]
{}^{FFM}_{0}D_{t}^{\alpha,\gamma} Z(\tilde{t}) = \mathcal{F}_4X, \\[6pt]
{}^{FFM}_{0}D_{t}^{\alpha,\gamma} Z_{a}(\tilde{t}) = \mathcal{F}_5X,
\end{cases}
\end{equation}
\begin{equation*}
X = (\tilde{t}, T, I, V, Z, Z_a)
\end{equation*}
where the nonlinear vector fields representing the transmission dynamics are defined as:
\begin{equation}
\begin{cases}
\mathcal{F}_1 = \lambda_T - (\mu_T + \chi V(\tilde{t}))T(\tilde{t}), \\[6pt]
\mathcal{F}_2 = \chi T(\tilde{t})V(\tilde{t}) - (\mu_I  + \alpha Z_{a}(\tilde{t}))I(\tilde{t}), \\[6pt]
\mathcal{F}_3 = \varepsilon_V \mu_I I(\tilde{t}) - \mu_V V(\tilde{t}), \\[6pt]
\mathcal{F}_4 = \lambda_Z - (\mu_Z  + \beta  I(\tilde{t}))Z(\tilde{t}), \\[6pt]
\mathcal{F}_5 = \beta Z(\tilde{t}) I(\tilde{t}) - \mu_{Z_a} Z_{a}(\tilde{t}).
\end{cases}
\end{equation}
The following results are obtained by using a fractal-fractional integral operator with a non-singular Mittag-Leffler law kernel, transforming the model into its corresponding Volterra integral form at the mesh point $t_{n+1}$:
\begin{equation}
\begin{aligned}
T(t_{n+1}) &= \frac{1-\alpha}{\mathbb{AB}(\alpha)} t_n^{\gamma-1} \mathcal{F}_1(t_n, P^n) + M \sum_{m=2}^{n} \int_{t_m}^{t_{m+1}} \mathcal{F}_1(\tau, P) \tau^{\gamma-1} (t_{n+1} - \tau)^{\alpha-1} d\tau, \\[6pt]
I(t_{n+1}) &= \frac{1-\alpha}{\mathbb{AB}(\alpha)} t_n^{\gamma-1} \mathcal{F}_2(t_n, P^n) + M \sum_{m=2}^{n} \int_{t_m}^{t_{m+1}} \mathcal{F}_2(\tau, P) \tau^{\gamma-1} (t_{n+1} - \tau)^{\alpha-1} d\tau, \\[6pt]
V(t_{n+1}) &= \frac{1-\alpha}{\mathbb{AB}(\alpha)} t_n^{\gamma-1} \mathcal{F}_3(t_n, P^n) + M \sum_{m=2}^{n} \int_{t_m}^{t_{m+1}} \mathcal{F}_3(\tau, P) \tau^{\gamma-1} (t_{n+1} - \tau)^{\alpha-1} d\tau, \\[6pt]
Z(t_{n+1}) &= \frac{1-\alpha}{\mathbb{AB}(\alpha)} t_n^{\gamma-1} \mathcal{F}_4(t_n, P^n) + M \sum_{m=2}^{n} \int_{t_m}^{t_{m+1}} \mathcal{F}_4(\tau, P) \tau^{\gamma-1} (t_{n+1} - \tau)^{\alpha-1} d\tau, \\[6pt]
Z_a(t_{n+1}) &= \frac{1-\alpha}{\mathbb{AB}(\alpha)} t_n^{\gamma-1} \mathcal{F}_5(t_n, P^n) + M \sum_{m=2}^{n} \int_{t_m}^{t_{m+1}} \mathcal{F}_5(\tau, P) \tau^{\gamma-1} (t_{n+1} - \tau)^{\alpha-1} d\tau,
\end{aligned}
\end{equation}
where
\begin{equation*}
    M = \frac{\alpha}{\mathbb{AB}(\alpha)\Gamma(\alpha)}.
\end{equation*}
$\mathbb{AB}(\alpha)$ is the normalization function of the Atangana-Baleanu kernel, and $P = \{T, I, V, Z, Z_a\}$ denotes the state vector layout. With the help of Newton polynomial, we approximate the combined continuous integrand function $\mathcal{G}(\tau, P) = \tau^{\gamma-1}\mathcal{F}_i(\tau, P)$ over the subinterval $[t_m, t_{m+1}]$ utilizing previous history keys $\{t_{m-2}, t_{m-1}, t_m\}$:
\begin{equation}
\begin{aligned}
\mathcal{G}(\tau, P) \simeq & \,\, \mathcal{G}(t_{m-2}, P^{m-2}) + \frac{1}{\Delta t} \left[ \mathcal{G}(t_{m-1}, P^{m-1}) - \mathcal{G}(t_{m-2}, P^{m-2}) \right] (\tau - t_{m-2}) \\
& + \frac{1}{2(\Delta t)^2} \left[ \mathcal{G}(t_m, P^m) - 2\mathcal{G}(t_{m-1}, P^{m-1}) + \mathcal{G}(t_{m-2}, P^{m-2}) \right] (\tau - t_{m-2})(\tau - t_{m-1}).
\end{aligned}
\end{equation}
For the integrals in the above formulation, we perform the following exact analytical integrations:
\begin{equation}
\begin{aligned}
\int_{t_m}^{t_{m+1}} (t_{n+1} - \tau)^{\alpha-1} d\tau &= \frac{(\Delta t)^\alpha}{\alpha} \left[ \Pi_k^n \right], \\[6pt]
\int_{t_m}^{t_{m+1}} (\tau - t_{m-2})(t_{n+1} - \tau)^{\alpha-1} d\tau &= \frac{(\Delta t)^{\alpha+1}}{\alpha(\alpha+1)} \left[ \Sigma_k^n \right], \\[6pt]
\int_{t_m}^{t_{m+1}} (\tau - t_{m-2})(\tau - t_{m-1})(t_{n+1} - \tau)^{\alpha-1} d\tau &= \frac{(\Delta t)^{\alpha+2}}{\alpha(\alpha+1)(\alpha+2)} \left[ \Omega_k^n \right],
\end{aligned}
\end{equation}
where the specialized algebraic weight constants tracking the historical trajectories are evaluated as:
\begin{equation}
\begin{aligned}
\Pi_k^n = & \, (n - m + 1)^\alpha - (n - m)^\alpha, \\
\Sigma_k^n = & \, (n - m + 1)^\alpha (n - m + 3 + 2\alpha) - (n - m)^\alpha (n - m + 3 + 3\alpha), \\
\Omega_k^n = & \, (n - m + 1)^\alpha \left\{ 2(n - m)^2 + (3\alpha + 10)(n - m) + 2\alpha^2 + 9\alpha + 12 \right\} \\
& - (n - m)^\alpha \left\{ 2(n - m)^2 + (5\alpha + 10)(n - m) + 6\alpha^2 + 18\alpha + 12 \right\}.
\end{aligned}
\end{equation}
Replacing the Newton polynomial terms back into the integrated architecture yields the complete multi-step numerical scheme for the healthy target CD4+ T-cell population $T(t_{n+1})$:
\begin{equation}
\begin{aligned}
T(t_{n+1}) = & \, \frac{1-\alpha}{\mathbb{AB}(\alpha)} t_n^{\gamma-1} \mathcal{F}_1(t_n, P^n) + \frac{(\Delta t)^\alpha}{\mathbb{AB}(\alpha)\Gamma(\alpha+1)} \sum_{m=2}^{n} t_{m-2}^{\gamma-1} \mathcal{F}_1(t_{m-2}, P^{m-2}) \times \left[ \Pi_k^n \right] \\
& + \frac{(\Delta t)^\alpha}{\mathbb{AB}(\alpha)\Gamma(\alpha+2)} \sum_{m=2}^{n} \left[ t_{m-1}^{\gamma-1} \mathcal{F}_1(t_{m-1}, P^{m-1}) - t_{m-2}^{\gamma-1} \mathcal{F}_1(t_{m-2}, P^{m-2}) \right] \times \left[ \Sigma_k^n \right] \\
& + \frac{(\Delta t)^\alpha}{2\mathbb{AB}(\alpha)\Gamma(\alpha+3)} \sum_{m=2}^{n} \left[ t_m^{\gamma-1} \mathcal{F}_1(t_m, P^m) - 2t_{m-1}^{\gamma-1} \mathcal{F}_1(t_{m-1}, P^{m-1}) + t_{m-2}^{\gamma-1} \mathcal{F}_1(t_{m-2}, P^{m-2}) \right] \times \left[ \Omega_k^n \right].
\end{aligned}
\end{equation}
By using the same process on the remaining structural compartments of the HIV dynamics, we complete the numerical scheme for computational simulations:
\begin{equation}
\begin{aligned}
I(t_{n+1}) = & \, \frac{1-\alpha}{\mathbb{AB}(\alpha)} t_n^{\gamma-1} \mathcal{F}_2(t_n, P^n) + \frac{(\Delta t)^\alpha}{\mathbb{AB}(\alpha)\Gamma(\alpha+1)} \sum_{m=2}^{n} t_{m-2}^{\gamma-1} \mathcal{F}_2(t_{m-2}, P^{m-2}) \times \left[ \Pi_k^n \right] \\
& + \frac{(\Delta t)^\alpha}{\mathbb{AB}(\alpha)\Gamma(\alpha+2)} \sum_{m=2}^{n} \left[ t_{m-1}^{\gamma-1} \mathcal{F}_2(t_{m-1}, P^{m-1}) - t_{m-2}^{\gamma-1} \mathcal{F}_2(t_{m-2}, P^{m-2}) \right] \times \left[ \Sigma_k^n \right] \\
& + \frac{(\Delta t)^\alpha}{2\mathbb{AB}(\alpha)\Gamma(\alpha+3)} \sum_{m=2}^{n} \left[ t_m^{\gamma-1} \mathcal{F}_2(t_m, P^m) - 2t_{m-1}^{\gamma-1} \mathcal{F}_2(t_{m-1}, P^{m-1}) + t_{m-2}^{\gamma-1} \mathcal{F}_2(t_{m-2}, P^{m-2}) \right] \times \left[ \Omega_k^n \right],
\end{aligned}
\end{equation}

\begin{thebibliography}{99}
\bibitem[7]{7} A. Atangana, Modelling the spread of COVID-19 with new fractal-fractional operators: can the lockdown save mankind before vaccination?, Chaos, Solitons \& Fractals, (2020) 136: 109860.
\end{thebibliography}
\end{document}
